\documentclass[11pt]{article}

\usepackage[a4paper,margin=1in]{geometry}
\usepackage{amsmath,amssymb}
\usepackage{booktabs}
\usepackage{graphicx}
\usepackage{hyperref}
\usepackage{enumitem}
\usepackage{xcolor}
\usepackage{listings}
\usepackage[T1]{fontenc}
\usepackage[utf8]{inputenc}

\title{BlackRay Deformed Kerr\newline Project Overview, Logic, and Future Outlook}
\author{Generated for the blackray-deformed-kerr repository}
\date{\today}

\begin{document}

\maketitle

\section*{Executive Summary}

This document presents a full project perspective for the \texttt{blackray-deformed-kerr} codebase. The goal of the project is to provide a reproducible, modular, and extensible toolkit for backward ray tracing in deformed Kerr spacetimes, with emphasis on Johannsen-type metric deformations, geodesic propagation, ISCO computation, disk-impact diagnostics, and local reflection-spectrum convolution. The code is organized around a compact C++17 core, Python bindings for simulation and analysis, and notebook-based visualization for scientific communication.

From a scientific and engineering point of view, the project has three defining characteristics:

\begin{enumerate}[leftmargin=1.5em]
    \item It separates the physics into reusable modules instead of keeping everything in one monolithic driver.
    \item It exposes a stable C ABI for Python-side automation and testing.
    \item It places explicit emphasis on scientific validation, including Kerr-limit checks, null-norm diagnostics, pathology screening, and multi-point comparisons.
\end{enumerate}

The repository is therefore not only a computational tool, but also a structured framework for future research, model comparison, and publication-ready figure generation.

\section{Project Scope and High-Level Purpose}

The program is designed to model photons propagating in a deformed Kerr spacetime and to estimate observables such as:

\begin{itemize}[leftmargin=1.5em]
    \item disk impact radii,
    \item redshift factors $g = E_{\mathrm{obs}}/E_{\mathrm{em}}$,
    \item emission angles,
    \item ISCO locations,
    \item null-norm conservation diagnostics,
    \item and local spectral convolution results.
\end{itemize}

The underlying metric is the Johannsen family, where the four deformation parameters are represented schematically as

\begin{equation}
A_1 = 1 + \frac{\alpha_{13}}{r^3}, \qquad
A_2 = 1 + \frac{\alpha_{22}}{r^2}, \qquad
A_5 = 1 + \frac{\alpha_{52}}{r^2}, \qquad
f = \frac{\epsilon_3}{r}.
\end{equation}

In the limit where all deformation parameters vanish, the code reduces to the standard Kerr metric. That limit is explicitly used as a validation target throughout the project.

\section{Repository Structure and Main Components}

The repository is structured into a small number of clear source areas:

\begin{verbatim}
blackray-deformed-kerr/
├── CMakeLists.txt
├── Makefile
├── README.md
├── requirements.txt
├── include/
│   ├── metric.hpp
│   ├── diffeqs.hpp
│   ├── isco.hpp
│   ├── raytracer.hpp
│   ├── convolver.hpp
│   └── blackray_c_api.h
├── src/
│   ├── metric.cpp
│   ├── diffeqs.cpp
│   ├── isco.cpp
│   ├── raytracer.cpp
│   ├── convolver.cpp
│   ├── blackray_c_api.cpp
│   ├── blackray_cli.cpp
│   └── ...
├── python/
│   ├── pyblackray.py
│   └── run_simulation.py
├── tests/
│   ├── test_kerr_limit.py
│   ├── test_null_norm.py
│   ├── test_pathology_screening.py
│   └── test_production_validation.py
├── notebooks/
│   ├── generate_demo_notebook.py
│   └── demo_deformed_kerr_analysis.ipynb
├── docs/
│   └── project_overview.tex
├── data/
├── build/
├── _legacy_blackray/
└── _legacy_raytransfer/
\end{verbatim}

The high-level decomposition is as follows:

\begin{itemize}[leftmargin=1.5em]
    \item \textbf{Metric logic:} \texttt{include/metric.hpp}, \texttt{src/metric.cpp}
    \item \textbf{Geodesic and integrator logic:} \texttt{include/diffeqs.hpp}, \texttt{src/diffeqs.cpp}
    \item \textbf{ISCO solver:} \texttt{include/isco.hpp}, \texttt{src/isco.cpp}
    \item \textbf{Ray tracing and disk intersections:} \texttt{include/raytracer.hpp}, \texttt{src/raytracer.cpp}
    \item \textbf{Convolution and spectrum processing:} \texttt{include/convolver.hpp}, \texttt{src/convolver.cpp}
    \item \textbf{C ABI bridge:} \texttt{include/blackray_c_api.h}, \texttt{src/blackray_c_api.cpp}
    \item \textbf{CLI and Python automation:} \texttt{src/blackray_cli.cpp}, \texttt{python/pyblackray.py}, \texttt{python/run_simulation.py}
    \item \textbf{Scientific validation:} \texttt{tests/}
    \item \textbf{Publication material:} \texttt{notebooks/}
\end{itemize}

\section{Core C and Python API Summary}

\begin{itemize}[leftmargin=1.5em]
    \item \textbf{C: blackray\_trace\_grid} \texttt{(x\_obs, y\_obs, count, observer\_radius, inclination, spin, alpha13, alpha22, alpha52, epsilon3, r\_isco, r\_out, r\_disk, g\_factor, cos\_theta\_e, status)} \textemdash{} traces a screen grid backward, returns disk-hit radii, redshift factors, emission cosines, and per-ray terminal status.
    \item \textbf{C: blackray\_max\_null\_norm} \texttt{(x\_obs, y\_obs, count, observer\_radius, inclination, spin, alpha13, alpha22, alpha52, epsilon3, step, steps, maximum\_norm)} \textemdash{} checks null-norm drift across a grid and reports the largest deviation.
    \item \textbf{C: blackray\_metric\_diagnostics} \texttt{(r, theta, spin, alpha13, alpha22, alpha52, epsilon3, determinant, g\_phiphi)} \textemdash{} returns metric-signature diagnostics for a given spacetime point.
    \item \textbf{Py: trace\_grid} \texttt{(x\_obs, y\_obs, observer\_radius=1000.0, inclination=1.0, spin=0.5, alpha13=0.0, alpha22=0.0, alpha52=0.0, epsilon3=0.0, r\_isco=0.0, r\_out=50.0, library=None)} \textemdash{} Python wrapper over the C API; returns \texttt{(records, status)} with disk radius, $g$, and $\cos\theta_{e}$ columns.
    \item \textbf{Py: max\_null\_norm} \texttt{(x\_obs, y\_obs, observer\_radius=1000.0, inclination=1.0, spin=0.5, alpha13=0.0, alpha22=0.0, alpha52=0.0, epsilon3=0.0, step=0.01, steps=100, library=None)} \textemdash{} Python convenience function for null-norm validation on a screen grid.
    \item \textbf{Py: metric\_diagnostics} \texttt{(r, theta, spin=0.5, alpha13=0.0, alpha22=0.0, alpha52=0.0, epsilon3=0.0, library=None)} \textemdash{} quick metric-level sanity check returning determinant and $g_{\phi\phi}$.
    \item \textbf{Py: find\_isco} \texttt{(spin, alpha13=0.0, alpha22=0.0, alpha52=0.0, epsilon3=0.0, library=None)} \textemdash{} analytic Kerr benchmark helper for the zero-deformation ISCO radius.
    \item \textbf{CLI: blackray\_cli} \texttt{(--spin, --inclination, --rstep, --pstep, --r-out, --output, --config, ...)} \textemdash{} command-line driver for screen-grid production runs and quick checks.
    \item \textbf{Python driver: run\_simulation.py} \texttt{(--mass, --spin, --alpha13, --alpha22, --alpha52, --epsilon3, --inclination, --rstep, --pstep, --r-out, --output, --spectrum-output, --plot)} \textemdash{} high-level batch run script that saves ray diagnostics and optional spectral products.
    \item \textbf{Python loader: pyblackray.py} \textemdash{} loads the compiled shared library via \texttt{BLACKRAY\_LIBRARY} or the local \texttt{build/} directory and configures the ctypes signatures automatically.
    \item \textbf{Notebook helper: generate\_demo\_notebook.py} \textemdash{} builds the publication-style notebook and figure set without requiring \texttt{nbformat}.
    \item \textbf{Validation scripts: test\_kerr\_limit.py, test\_null\_norm.py, test\_pathology\_screening.py, test\_production\_validation.py} \textemdash{} verify the Kerr limit, null-norm stability, pathology handling, and production-grade comparisons.
\end{itemize}

\section{Core Logic by Module}

\subsection{1. Metric Module}

The metric module is the foundation of the project. It defines the Johannsen metric components, stores the compact $t$-$\phi$ block needed by the geodesic solver, and computes derivatives with respect to radius and polar angle.

\textbf{Key files}

\begin{itemize}[leftmargin=1.5em]
    \item \texttt{include/metric.hpp}
    \item \texttt{src/metric.cpp}
\end{itemize}

\textbf{Main responsibilities}

\begin{itemize}[leftmargin=1.5em]
    \item Validation of input parameters.
    \item Computation of the full metric tensor and compact metric components.
    \item Derivative calculations for radial and polar dependence.
    \item Inversion of the $t$-$\phi$ block for use in geodesic equations.
\end{itemize}

This module is the first place where the code distinguishes itself from a purely black-box legacy driver. Rather than hard-coding a single geometry, it provides a reusable metric interface that other modules can call consistently.

\subsection{2. Geodesic and Differential-Equation Module}

The differential-equation module is responsible for the first-order equations governing null geodesics and the numerical stepping strategy used to trace them. It handles both the actual photon propagation and the diagnostic checks that confirm the model remains physically meaningful.

\textbf{Key files}

\begin{itemize}[leftmargin=1.5em]
    \item \texttt{include/diffeqs.hpp}
    \item \texttt{src/diffeqs.cpp}
\end{itemize}

\textbf{Main responsibilities}

\begin{itemize}[leftmargin=1.5em]
    \item Definition of the geodesic state and integrator options.
    \item Evaluation of Christoffel symbols and the geodesic ODE system.
    \item Adaptive RK4-type stepping with protective checks.
    \item Computation of the null norm of the geodesic state.
\end{itemize}

In practical terms, this module determines whether a traced photon stays on a valid null trajectory, whether a step is too large, and whether the numerical evolution remains inside the physical domain.

\subsection{3. ISCO Module}

The ISCO module computes circular orbits and identifies the marginally stable orbit, which is the natural inner edge for a thin disk model in many ray-tracing applications.

\textbf{Key files}

\begin{itemize}[leftmargin=1.5em]
    \item \texttt{include/isco.hpp}
    \item \texttt{src/isco.cpp}
\end{itemize}

\textbf{Main responsibilities}

\begin{itemize}[leftmargin=1.5em]
    \item Compute Keplerian angular velocity $\Omega$.
    \item Infer $u^t$ for circular equatorial motion.
    \item Scan the radial domain for the ISCO root.
    \item Report auxiliary diagnostics such as energy, angular momentum, and vertical-instability warnings.
\end{itemize}

This part is especially important because the ray tracer relies on a physically meaningful inner radius for disk impacts. In the current implementation, the ISCO is interpreted as the inner edge of the disk in the standard ray-tracing workflow.

\subsection{4. Ray-Tracing Module}

The ray tracer carries the main observational logic: it takes an observer-screen coordinate, builds the local camera tetrad, traces photons backward, and records whether each ray reaches the disk, escapes, is halted by the horizon, or fails numerically.

\textbf{Key files}

\begin{itemize}[leftmargin=1.5em]
    \item \texttt{include/raytracer.hpp}
    \item \texttt{src/raytracer.cpp}
\end{itemize}

\textbf{Main responsibilities}

\begin{itemize}[leftmargin=1.5em]
    \item Camera geometry and observer setup.
    \item Construction of local tetrads.
    \item Backward integration from the observer to the disk or escape region.
    \item Disk-hit bookkeeping, including radial coordinate, redshift, and emission angle.
    \item Safe termination states for horizon, escape, and failure.
\end{itemize}

This is the central module that turns the abstract metric and geodesic physics into observable quantities for imaging and spectroscopy.

\subsection{5. Convolver Module}

The convolver module performs local spectral processing, including broadening and interpolation logic that mimics a simplified reflection-spectrum response.

\textbf{Key files}

\begin{itemize}[leftmargin=1.5em]
    \item \texttt{include/convolver.hpp}
    \item \texttt{src/convolver.cpp}
\end{itemize}

\textbf{Main responsibilities}

\begin{itemize}[leftmargin=1.5em]
    \item Local reflection-spectrum interpolation.
    \item Broadening of the emitted line profile.
    \item Generation of disk-averaged spectral diagnostics.
\end{itemize}

This module connects the geometry and propagation results to a more astrophysical interpretation of the emission spectrum.

\subsection{6. C ABI and CLI Layer}

The project exposes a C-language ABI for stable interoperability with Python and external tools.

\textbf{Key files}

\begin{itemize}[leftmargin=1.5em]
    \item \texttt{include/blackray_c_api.h}
    \item \texttt{src/blackray_c_api.cpp}
    \item \texttt{src/blackray_cli.cpp}
\end{itemize}

\textbf{Why this matters}

The C ABI provides a clean boundary between the C++ core and the external ecosystem. This is important because it allows the library to be loaded by Python via \texttt{ctypes}, while also keeping the core implementation separated from the scripting layer.

\subsection{7. Python Wrapper and Simulation Script}

The Python layer is intentionally lightweight and acts as a bridge between the compiled core and analysis workflows.

\textbf{Key files}

\begin{itemize}[leftmargin=1.5em]
    \item \texttt{python/pyblackray.py}
    \item \texttt{python/run_simulation.py}
\end{itemize}

\textbf{Responsibilities}

\begin{itemize}[leftmargin=1.5em]
    \item Load the shared library.
    \item Expose the core API to Python.
    \item Run batch screen simulations.
    \item Save ray tables and optional spectral outputs.
\end{itemize}

This design makes it possible to integrate the engine into notebooks, test scripts, and automated scientific pipelines without reimplementing the physics in Python.

\subsection{8. Tests and Validation}

The validation suite is one of the strongest parts of the current project structure.

\textbf{Key files}

\begin{itemize}[leftmargin=1.5em]
    \item \texttt{tests/test_kerr_limit.py}
    \item \texttt{tests/test_null_norm.py}
    \item \texttt{tests/test_pathology_screening.py}
    \item \texttt{tests/test_production_validation.py}
\end{itemize}

The tests cover:

\begin{itemize}[leftmargin=1.5em]
    \item the zero-deformation Kerr limit,
    \item null geodesic conservation,
    \item pathology screening in unphysical regions,
    \item and production-style validation across multiple spins and radii.
\end{itemize}

This is exactly the sort of validation needed before scientific claims are made about performance or model correctness.

\subsection{9. Notebook and Publication Materials}

The notebook generator creates a direct JSON-based notebook artifact without requiring \texttt{nbformat}.

\textbf{Key files}

\begin{itemize}[leftmargin=1.5em]
    \item \texttt{notebooks/generate_demo_notebook.py}
    \item \texttt{notebooks/demo_deformed_kerr_analysis.ipynb}
\end{itemize}

This part of the project is important for two reasons:

\begin{enumerate}[leftmargin=1.5em]
    \item It produces a reproducible research demo.
    \item It provides a clear bridge from code to figures for papers and presentations.
\end{enumerate}

\section{Key Logic Differences from the Legacy Workflow}

The current project is not simply a line-for-line reimplementation of the legacy ABHModels workflow. It deliberately separates concerns and upgrades the engineering model in several important ways.

\subsection{1. Explicit module boundaries}

In the legacy code, the physics was often intermixed across drivers and helper routines. In the new code, the metric, the geodesics, the ISCO solver, the ray tracer, and the spectrum logic are all distinct software components. This improves clarity, validation, and reuse.

\subsection{2. Stable scientific checks}

The new implementation includes explicit checks for:

\begin{itemize}[leftmargin=1.5em]
    \item invalid inputs,
    \item singular or ill-conditioned metrics,
    \item unphysical circular orbits,
    \item and null-norm drift.
\end{itemize}

These checks are scientifically useful because they prevent silent numerical pathologies from being mistaken for physical results.

\subsection{3. Python accessibility}

The project exposes the main numerical functionality through Python, which is important for tests, notebook analysis, and reproducible scientific workflows. The legacy code relied on a more fragmented set of tools and driver scripts.

\subsection{4. Scientific validation first}

A major philosophical shift is that the project now treats the Kerr limit and analytical comparisons as a required scientific baseline instead of only a smoke test. This is a stronger and more publication-friendly way to build confidence in the implementation.

\subsection{5. Publication-oriented output}

The notebook generator and the figure-producing cells help convert simulation results into publication-ready figures quickly. This reduces the friction between numerical experimentation and scholarly communication.

\section{A Full Project Perspective (FPP)}

From a full project-management perspective, the repository sits at the intersection of three domains:

\begin{enumerate}[leftmargin=1.5em]
    \item \textbf{Physics modeling}: deformed Kerr metrics, geodesic propagation, ISCO physics, and disk observables.
    \item \textbf{Software systems engineering}: modular C++ core, C ABI, Python wrappers, and testing infrastructure.
    \item \textbf{Scientific communication}: notebooks, plots, validation summaries, and reproducible analysis.
\end{enumerate}

In that sense, the project is not merely a solver. It is a small research platform for relativistic astrophysics workflows. Its strongest value lies in the combination of mathematical clarity, code modularity, and validation discipline.

\section{Future Outlook and Roadmap}

Several useful directions are now natural extensions of the present codebase.

\subsection{1. More complete spectral and transfer-function comparisons}

The current notebook already includes several families of comparison plots. The next natural step is to formalize a broader figure set covering:

\begin{itemize}[leftmargin=1.5em]
    \item transfer-function variations with spin,
    \item transfer-function variations with inclination,
    \item deformation-parameter sensitivity maps,
    \item incident-flux profiles and line-shape trends,
    \item and explicit comparison to legacy RELXILL or RELXILL	extsubscript{NK}-style contexts.
\end{itemize}

\subsection{2. Parameter sweeps and grid-based studies}

A publication-quality workflow would benefit from automated parameter sweeps over:

\begin{itemize}[leftmargin=1.5em]
    \item $a^*$,
    \item $\alpha_{13}$, $\alpha_{22}$, $\epsilon_3$, $\alpha_{52}$,
    \item inclination,
    \item and lamppost height.
\end{itemize}

These sweeps could be organized into structured data products and saved as machine-readable tables for later fitting and plotting.

\subsection{3. Tighter integration with external spectral model libraries}

The current code already has a convolver module, but the strongest long-term extension would be direct integration with established reflection and disk models. This would make it easier to compare the deformed-Kerr framework against external spectral fitting tools.

\subsection{4. More rigorous numerical convergence studies}

A further step would be to quantify convergence across screen resolution, integration tolerances, and radial sampling for both the ray-tracing and ISCO modules. This would make the results more credible from a numerical-analysis perspective.

\subsection{5. Better reproducibility and packaging}

Improving packaging and release metadata would be valuable for external users. This includes a clear publication-ready versioning strategy, a proper license review, and a defined citation workflow.

\section{Representative Parameter Summary}

The following summary provides a compact overview of the simulation/fitting parameters that are relevant to the current project workflow.

\begin{table}[h]
\centering
\begin{tabular}{ll}
\toprule
\textbf{Quantity} & \textbf{Representative value / range} \\
\midrule
Dimensionless spin $a^*$ & $0.98$ baseline; also used in sweeps such as $[-0.5,0.5,0.98]$ \\
Inclination $\iota$ & $30^\circ$ baseline; also used in sweeps such as $[10^\circ,30^\circ,50^\circ,70^\circ]$ \\
Lamppost height $h$ & $2$ and $10$ in representative flux-profile comparisons \\
Deformation parameters & $\alpha_{13},\alpha_{22},\epsilon_3,\alpha_{52}$ in sensitivity sweeps \\
Photon index $\Gamma$ & $2.0$ \\
Emissivity index $q$ & $3$ \\
Cosmological redshift $z$ & $0$ \\
Ionization parameter $\log\xi$ & $3.1$ \\
Iron abundance $A_{\mathrm{Fe}}$ & $1.0$ (solar) \\
Outer disk radius $r_{\mathrm{out}}$ & $30$ \\
Observer radius $r_{\mathrm{obs}}$ & $100$ \\
Screen grid size & $8\times 8$ or $16\times 16$ depending on the figure or test \\
\bottomrule
\end{tabular}
\caption{Representative parameter values used in the current project workflow.}
\end{table}

\section{Concluding Remarks}

The repository already embodies a coherent scientific project: it combines a physical spacetime model, numerical propagation, a clean software architecture, and publication-oriented diagnostics. The present code is stronger than a simple legacy port because it explicitly validates the Kerr limit, separates the mathematics from the user-facing layers, and provides a reproducible path from simulation to figures.

In short, the project is best understood as a modular platform for deformed-Kerr ray tracing and spectroscopy that is already capable of supporting both research and communication workflows, while remaining open to continued expansion in parameter sweeps, model comparison, and production-level validation.

\end{document}
